\documentclass[11pt,a4paper]{article}

\usepackage[margin=1in]{geometry}
\usepackage{amsmath, amssymb, amsthm}
\usepackage{mathtools}
\usepackage{bm}
\usepackage{enumitem}
\usepackage{microtype}
\usepackage[T1]{fontenc}
\usepackage{lmodern}
\usepackage[dvipsnames]{xcolor}
\usepackage{hyperref}

\hypersetup{
  colorlinks=true,
  linkcolor=blue!50!black,
  urlcolor=blue!50!black,
  citecolor=blue!50!black
}

\newtheorem{theorem}{Theorem}[section]
\newtheorem{lemma}[theorem]{Lemma}
\newtheorem{definition}[theorem]{Definition}

\title{U4: The Forced Structure of the Non-Relativistic Hamiltonian}

\author{Allen Proxmire}

\date{April 2026}

\begin{document}

\maketitle

\begin{abstract}
We present the structural derivation of the non-relativistic single-particle Hamiltonian $\hat{H} = \hbar^2 |\hat{p}|^2 / (2m) + V(\hat{x})$ from the Event-Density (ED) primitive stack and the now-\textsc{forced} upstream items U1, U2, U5, and Theorem~10. Operating under Framing~1 (U4 conditional on the existence of a Hermitian Hamiltonian generator, which is sibling \textsc{candidate} U3's content), we decompose U4 into five sub-features --- channel--momentum identification (F1), plane-wave eigenfunctions (F2), quadratic-in-$|\hat{p}|^2$ kinetic structure (F3), specific coefficient $1/(2m)$ (F4), and kinetic + potential decomposition (F5) --- and derive each from primitive-level inputs plus standard mathematical infrastructure. F1 follows from Stone's theorem applied to translation symmetry on the U2 Hilbert space, transferring directly from the U5 derivation; F2 is mathematical; F5 follows from translation invariance + locality. F3 is established via four constraints (translation, rotation, analyticity, non-relativistic-limit), with all candidate alternatives --- linear-in-$|\hat{p}|$, anisotropic, higher-even-power, full relativistic dispersion --- explicitly dismissed within the non-relativistic scope. F4 is the methodologically distinctive content: \textbf{the factor of 2 in $1/(2m)$ is forced by integration of the Galilean commutator condition $[\hat{H}, \hat{K}_i] = -i\hbar \hat{p}_i$ against the boost generator $\hat{K} = m\hat{x} - t\hat{p}$, with the factor of 2 emerging as the integration Jacobian rather than as a convention borrowed from classical mechanics.} All seven falsifiers are dispatched within the stated scope conditions. The U4 verdict is \textbf{\textsc{forced} with form-\textsc{forced}-value-\textsc{inherited} framing}: the kinetic-energy form is uniquely derived from primitives + symmetry + scope; the mass parameter $m$ is inherited per Arc~M's ``form \textsc{forced}, values \textsc{inherited}'' methodology, and $\hbar$ is inherited via the dimensional-atlas Madelung anchoring. The arc introduces no new structural \textsc{candidate}s; the conditionalities are all inherited. The active upstream-\textsc{candidate} inventory of the QM-emergence Phase-1 program reduces from $\{$U3, U4$\}$ + continuum gauge to \textbf{$\{$U3$\}$ + continuum gauge} --- one upstream item remains to close the entire Phase-1 structural-foundations program. With U4 closed, the Phase-1 theorem inventory stands at 15 forced structural theorems; Schr\"odinger emergence is now \textsc{forced} conditional on U3 alone.
\end{abstract}

\section{Introduction}

The Schr\"odinger equation is one of the most directly tested structures in physics. Its non-relativistic single-particle form, $i\hbar \partial_t \psi = \hat{H} \psi$ with $\hat{H} = -\hbar^2 \nabla^2 / (2m) + V(\hat{x})$, has been used in calculations of atomic spectra, molecular bonding, semiconductor physics, neutron interferometry, and countless other empirical settings for nearly a century. The Hamiltonian's specific form --- the squared-momentum kinetic energy with the specific factor $1/(2m)$, plus the additive position-dependent potential --- is the structural anchor that connects the abstract evolution equation to physical observables.

In standard quantum mechanics, this Hamiltonian form is \emph{imported} rather than derived. The kinetic-energy expression $T = p^2/(2m)$ is taken from classical mechanics, where it is the kinetic energy of a non-relativistic point particle; the operator-replacement $p \to \hat{p}$ promotes it to operator status. The factor of 2 in $1/(2m)$ is borrowed from the classical formula without further justification.

The Event-Density (ED) program asks whether this structural form can be derived rather than postulated. The QM-emergence Phase-1 framework~\cite{phase1} presents non-relativistic single-particle quantum mechanics as a structural consequence of a primitive ontology. Five arcs of the structural-foundations cycle (born\_gleason~\cite{borngleason}, U2-Discrete~\cite{u2}, U2-Continuum~\cite{u2}, U5~\cite{u5}, U1~\cite{u1}) have promoted the participation-measure carrier, its inner-product structure, the Born rule, the adjacency-band Fourier-conjugate partition, and the participation-measure construction itself from \textsc{candidate} to \textsc{forced} status. With those five arcs closed, two upstream commitments remained: U3 (the participation-measure evolution equation, supplying existence + linear-first-order structure of the time-evolution generator) and U4 (the specific Hamiltonian form, supplying the dynamical content of the kinetic energy). This paper presents the U4 arc and its verdict.

The U4 arc is the sixth and most delicate of the structural-foundations cycle. It is the first arc whose verdict is naturally framed as conditional on a \emph{sibling \textsc{candidate}} (U3) --- methodologically novel for the cycle --- and it is the first arc whose load-bearing content concerns \emph{dynamical structure} (the time-evolution generator's form) rather than kinematic structure. The arc's substantive structural finding is the derivation of the factor of 2 in $1/(2m)$ as an \emph{integration Jacobian} of the Galilean Lie algebra commutator condition, rather than as a convention borrowed from classical mechanics. The arc closes \textsc{forced} at the form level, with the mass parameter $m$ and the constant $\hbar$ inherited per established ED-program methodology.

The paper is organized as follows. Section~\ref{sec:question} states the U4 question precisely. Section~\ref{sec:inputs} inventories the structural inputs available under Framing~1. Section~\ref{sec:decomposition} decomposes U4 into five sub-features. Section~\ref{sec:falsifiers} introduces the seven falsifiers. Sections~\ref{sec:F1}--\ref{sec:F4} derive the five sub-features in turn, with Sections~\ref{sec:F3} and~\ref{sec:F4} carrying the load-bearing content. Section~\ref{sec:falsifier-table} consolidates falsifier dispatches. Section~\ref{sec:conditionality} documents the conditionality ledger. Section~\ref{sec:final-theorem} states the final U4 theorem. Section~\ref{sec:downstream} documents the downstream cascade. Section~\ref{sec:discussion} discusses the result's significance.

\section{Statement of the U4 Question}\label{sec:question}

\subsection{The target form}

The QM-emergence Phase-1 framework's Step 2 (Schr\"odinger emergence) targets the standard non-relativistic single-particle Hamiltonian:
\begin{equation}
\hat{H} = \frac{\hbar^2 |\hat{p}|^2}{2m} + V(\hat{x}) \qquad (\text{U4})
\label{eq:U4target}
\end{equation}
with the kinetic operator $T(\hat{p}) = \hbar^2 |\hat{p}|^2 / (2m)$ depending only on momentum and the potential operator $V(\hat{x})$ depending only on position. The participation-measure framework's evolution equation (U3) is
\[
i\hbar \partial_t P_K = \hat{H} P_K + \sum_{K'} V_{KK'} P_{K'}.
\]
U4 specifies the \emph{form} of $\hat{H}$ entering this equation; U3 supplies the \emph{existence + linear-first-order structure} of the equation itself.

\subsection{The five sub-features}

\begin{itemize}
  \item \textbf{F1 (Channel--momentum identification).} In the thin/continuum limit, the channel index $K$ admits a continuous momentum-space coordinate $k$ via Stone's theorem on the participation graph's translation group.
  \item \textbf{F2 (Plane-wave eigenfunctions).} The eigenfunctions of the translation generator $\hat{p}$ in position representation are plane waves $\langle x | k \rangle = (2\pi\hbar)^{-d/2} e^{i k x / \hbar}$.
  \item \textbf{F3 (Quadratic-in-$|\hat{p}|^2$ kinetic structure).} $T(\hat{p}) \propto |\hat{p}|^2$, with linear-in-$|\hat{p}|$, higher-even-power, anisotropic, and full relativistic alternatives all excluded in the non-relativistic regime.
  \item \textbf{F4 (Specific coefficient $1/(2m)$).} $T(\hat{p}) = \hbar^2 |\hat{p}|^2 / (2m)$ --- the factor of 2 in the denominator is structural, not conventional.
  \item \textbf{F5 (Kinetic + potential decomposition).} $\hat{H} = T(\hat{p}) + V(\hat{x})$ with no cross terms, in the gauge-coupling-free scope.
\end{itemize}

\subsection{The seven falsifiers}

\begin{itemize}
  \item \textbf{Fal-1.} No momentum-basis identification.
  \item \textbf{Fal-2.} Linear-in-$|\hat{p}|$ kinetic term (e.g., $T = c |\hat{p}|$, photon-like dispersion).
  \item \textbf{Fal-3.} Higher-even-power terms (e.g., $T \supset c |\hat{p}|^4$ in the non-relativistic regime).
  \item \textbf{Fal-4.} Coefficient $\neq 1/(2m)$ --- wrong factor.
  \item \textbf{Fal-5.} Anisotropic kinetic energy (direction-dependent).
  \item \textbf{Fal-6.} Position-momentum cross terms in $\hat{H}$.
  \item \textbf{Fal-7.} Full relativistic dispersion in non-relativistic regime.
\end{itemize}

\subsection{Framing 1: U4 conditional on sibling \textsc{candidate} U3}

The U4 arc adopts \textbf{Framing~1}: U4 is derived \emph{given} that a Hermitian Hamiltonian generator $\hat{H}$ exists on the participation-measure Hilbert space. The existence question is U3's content; U4 specifies the form. This framing is methodologically novel within the structural-foundations cycle.

Under Framing~1, U3 cannot be invoked as a \emph{derivation input} for U4 (that would be circular), but the existence of a Hermitian generator can be assumed as the starting point. The verdict is then ``U4 \textsc{forced} conditional on U3''; promoting U3 in a subsequent arc would make U4 fully unconditional.

\section{Structural Inputs}\label{sec:inputs}

\subsection{Primitives}

\begin{itemize}
  \item \textbf{Primitive 03 (Participation):} graph homogeneity, supplying translation symmetry kinematically.
  \item \textbf{Primitive 06 (ED Gradient):} $\nabla \rho$ supplies a spatial direction with no preferred origin.
  \item \textbf{Primitive 09 (Tension Polarity):} $U(1)$-valued polarity phase entering the participation-measure construction via U1.
  \item \textbf{Primitive 13 (Relational Timing):} time-axis structure, used in F4 to combine with spatial translation symmetry into the full Galilean group.
\end{itemize}

\subsection{Inherited \textsc{forced} upstream items}

The U4 arc has access to all five \textsc{forced} items from prior arcs of the structural-foundations cycle: U1~\cite{u1}, U2-Discrete and U2-Continuum~\cite{u2}, U5~\cite{u5}, Theorem~10~\cite{borngleason}. Particularly, the Stone's-theorem derivation in U5 transfers directly to F1 of the present arc.

\subsection{Mathematical infrastructure}

\begin{itemize}
  \item Stone's theorem on one-parameter unitary groups (Stone 1932)~\cite{stone}.
  \item Standard functional calculus on self-adjoint operators.
  \item Galilean Lie algebra and central extension (Bargmann 1954)~\cite{bargmann}.
  \item Standard Fourier transform on $L^2(\mathbb{R}^d)$.
\end{itemize}

\subsection{Scope conditions}

Two scope conditions inherited from Phase-1 Step 2's framing:
\begin{itemize}
  \item \textbf{Non-relativistic single-particle scope.} $|\hat{p}|^2 / (mc)^2 \ll 1$.
  \item \textbf{Gauge-coupling-free scope.} Magnetic vector potentials would couple position and momentum via $T(\hat{p} - eA(\hat{x})/c)$; F5 holds in the absence of such gauge couplings.
\end{itemize}

\subsection{Inherited values}

\begin{itemize}
  \item \textbf{Mass $m$:} per Arc~M's ``form \textsc{forced}, values \textsc{inherited}'' methodology~\cite{arcm}.
  \item \textbf{$\hbar$:} inherited via the dimensional-atlas Madelung anchoring.
\end{itemize}

\subsection{Forbidden inputs (under Framing 1)}

The following inputs are explicitly forbidden as derivation inputs to U4: U3 (sibling \textsc{candidate}; may be invoked only as conditional); the classical Hamiltonian as a postulate; the Schr\"odinger equation in textbook form; the relativistic dispersion as a derivation input.

\section{Sub-Feature Decomposition}\label{sec:decomposition}

\begin{center}
\begin{tabular}{p{3.0cm}p{3.5cm}p{6.0cm}p{2.5cm}}
\hline
Sub-feature & Status & Structural inputs & Falsifiers \\
\hline
F1 & \textsc{forced} & Primitives 03 + 06 + U2 + Stone's theorem; U3-independent & Fal-1 \\
F2 & \textsc{forced} (automatic given F1) & F1 + standard math & --- \\
F3 & \textsc{forced} & C1 translation + C2 rotation + C3 analyticity + C4 non-rel limit & Fal-2, Fal-3, Fal-5, Fal-7 \\
F4 & \textsc{forced} form (factor of 2 derived); $m$, $\hbar$ \textsc{inherited} & Primitives 03 + 13 + non-rel scope + U2 unitarity + Galilean Lie algebra integration & Fal-4 \\
F5 & \textsc{forced} (gauge-coupling-free scope) & Translation invariance + locality & Fal-6 \\
\hline
\end{tabular}
\end{center}

\section{Falsifier Inventory}\label{sec:falsifiers}

Each falsifier targets a specific physical alternative:
\begin{itemize}
  \item \textbf{Fal-1:} Channel index has no momentum-space interpretation. Tested by Stone's theorem on the U2 Hilbert space.
  \item \textbf{Fal-2:} Massless photon-like dispersion in a regime supposed to describe massive particles. Tested by C2 + C3 + C4.
  \item \textbf{Fal-3:} Relativistic corrections in non-relativistic regime. Tested by C4 scope condition.
  \item \textbf{Fal-4:} Wrong proportionality between kinetic energy and momentum-squared. Tested by Galilean Lie algebra integration.
  \item \textbf{Fal-5:} Anisotropic kinetic energy. Tested by C2 rotation invariance.
  \item \textbf{Fal-6:} Position-momentum cross terms. Tested by translation invariance + locality.
  \item \textbf{Fal-7:} Full relativistic dispersion. Tested by C4 non-relativistic-limit.
\end{itemize}

\section{Derivation of F1 (Momentum Identification)}\label{sec:F1}

\subsection{The claim}

In the thin/continuum limit, the channel index $K$ admits a continuous momentum-space coordinate $k$ via Stone's theorem on the participation graph's translation group, with the basis-change between channel-representation and momentum-representation the standard Fourier transform.

\subsection{Translation symmetry as a kinematic property}

Define spatial translation $T_a$ acting on the participation-measure space by $(T_a P)_K(u) := P_K(u - a)$. Primitive~03's homogeneity guarantees that the participation relation does not privilege any vertex. The set $\{T_a : a \in \mathbb{R}^d\}$ forms a $d$-parameter abelian Lie group.

\subsection{Translation acts unitarily on the U2 Hilbert space}

For any $P, Q \in \mathcal{H}$ and any translation $a$:
\[
\langle T_a P \mid T_a Q \rangle = \sum_K \int du \; P_K^*(u - a) \, Q_K(u - a) = \sum_K \int du' \; P_K^*(u') \, Q_K(u') = \langle P \mid Q \rangle
\]
under change of variables $u' = u - a$. The translation group preserves the U2 inner product. Continuity follows from continuity of the participation-measure field. Therefore $\{T_a\}$ is a strongly continuous one-parameter unitary group on $\mathcal{H}$.

\subsection{Stone's theorem identifies the translation generator}

\begin{theorem}[Stone 1932~\cite{stone}]
Every strongly continuous one-parameter unitary group on a Hilbert space has a unique self-adjoint generator.
\end{theorem}

Applied to the spatial translation group:
\begin{equation}
T_a = \exp\!\left( \frac{i \hat{p} \cdot a}{\hbar} \right). \label{eq:Ta}
\end{equation}
The eigenstates $|k\rangle$ of $\hat{p}$ form a complete orthonormal basis on $\mathcal{H}$. \textbf{In the thin/continuum limit, the channel index $K$ is identified with the momentum eigenvalue $k$:}
\begin{equation}
K \leftrightarrow k. \label{eq:Kk}
\end{equation}

\subsection{Explicit U3-independence}

The derivation invokes Primitives 03 + 06, U2, and Stone's theorem. \textbf{No time-evolution operator, no Hamiltonian, no Schr\"odinger equation appears in the argument.} Stone's theorem identifies $\hat{p}$ kinematically on the U2 Hilbert space.

\begin{theorem}[F1]
Let $\mathcal{H}$ be the U2 Hilbert space of participation measures. Let $\{T_a : a \in \mathbb{R}^d\}$ be the spatial translation group acting on $\mathcal{H}$. Then $\{T_a\}$ is a strongly continuous one-parameter unitary group, and there exists a unique self-adjoint operator $\hat{p}$ such that $T_a = \exp(i \hat{p} \cdot a / \hbar)$. The eigenstates of $\hat{p}$ form a complete orthonormal basis on $\mathcal{H}$, and in the thin/continuum limit the channel index $K$ is identified with the eigenvalue $k$ of $\hat{p}$.
\end{theorem}

\textbf{Falsifier Fal-1 dispatched.}

\section{Derivation of F2 (Plane-Wave Eigenfunctions)}\label{sec:F2}

\subsection{Position-representation form of $\hat{p}$}

Given F1's identification of $\hat{p}$, the position-representation form is determined by:
\[
T_a \psi(x) = \psi(x - a) = \exp\!\left( -a \cdot \nabla \right) \psi(x).
\]
Comparing with~\eqref{eq:Ta}:
\begin{equation}
\hat{p} = -i \hbar \nabla. \label{eq:p-position}
\end{equation}

\subsection{Eigenfunctions}

The eigenvalue equation $\hat{p} \, |k\rangle = k \, |k\rangle$ becomes, in position representation, $-i \hbar \nabla \langle x | k \rangle = k \langle x | k \rangle$, with solution
\begin{equation}
\langle x | k \rangle = (2\pi\hbar)^{-d/2} \cdot e^{i k \cdot x / \hbar}. \label{eq:planewave}
\end{equation}
The normalization $(2\pi\hbar)^{-d/2}$ is fixed by the orthonormality $\langle k | k' \rangle = \delta(k - k')$.

\begin{theorem}[F2]
In the position representation of the U2 Hilbert space, the translation generator $\hat{p}$ takes the form $\hat{p} = -i \hbar \nabla$. The eigenfunctions of $\hat{p}$ are plane waves $\langle x | k \rangle = (2\pi\hbar)^{-d/2} \cdot e^{i k \cdot x / \hbar}$.
\end{theorem}

F2 introduces no new structural commitment beyond F1.

\section{Derivation of F5 (Kinetic + Potential Decomposition)}\label{sec:F5}

\subsection{Translation invariance forces the free part to depend only on $\hat{p}$}

In the absence of an external potential, $\hat{H}_\mathrm{free}$ commutes with translations: the participation graph itself is translation-invariant (Primitive~03), and without external coupling, $\hat{H}_\mathrm{free}$ inherits this invariance. Any operator commuting with $\hat{p}$ depends on $\hat{p}$ alone via standard functional calculus. Therefore
\begin{equation}
\hat{H}_\mathrm{free} = T(\hat{p}). \label{eq:Hfree}
\end{equation}

\subsection{Locality forces external potentials to depend only on $\hat{x}$}

Locality is a primitive-level structural property of the participation graph: edges connect adjacent vertices, and external influences couple at specific locations. The local action of a potential takes the form $(V \psi)(x) = V(x) \psi(x)$ --- multiplication by a position-dependent function, depending only on $\hat{x}$.

A potential depending on $\hat{p}$ would be non-local in position representation and is structurally excluded from the standard non-relativistic-single-particle scope.

\subsection{Cross terms forbidden}

Any cross term $\hat{p} \hat{x}$ or $\hat{x} \hat{p}$ violates either translation invariance (free part) or locality (potential part).

\subsection{Scope caveat}

In the presence of magnetic vector potentials $A(\hat{x})$, the kinetic operator becomes $T(\hat{p} - e A(\hat{x})/c)$. The scope condition for U4's F5 is non-relativistic single-particle Schr\"odinger emergence \emph{without} gauge couplings.

\begin{theorem}[F5]
Within the gauge-coupling-free non-relativistic single-particle scope of Phase-1 Step~2, the Hamiltonian on the U2 Hilbert space admits the unique additive decomposition $\hat{H} = T(\hat{p}) + V(\hat{x})$, with the kinetic part depending only on $\hat{p}$ and the potential part depending only on $\hat{x}$. Cross terms are forbidden.
\end{theorem}

\textbf{Falsifier Fal-6 dispatched.}

\section{Derivation of F3 (Quadratic Kinetic Structure)}\label{sec:F3}

\subsection{The four constraints}

\begin{itemize}
  \item \textbf{C1 Translation invariance.} $T = T(\hat{p})$ alone (already established).
  \item \textbf{C2 Rotation invariance.} The participation graph supports rotations as a kinematic symmetry.
  \item \textbf{C3 Analyticity.} $T(\hat{p}) = f(|\hat{p}|^2)$ is analytic in $|\hat{p}|^2$ at $\hat{p} = 0$.
  \item \textbf{C4 Non-relativistic-limit scope.} $|\hat{p}|^2 / (mc)^2 \ll 1$.
\end{itemize}

\subsection{C1 + C2 reduce to $f(|\hat{p}|^2)$}

By C1, $T = T(\hat{p})$. By C2, $T$ is invariant under rotations of $\hat{p}$. The unique rotational invariant of $\hat{p}$ (up to powers) is $|\hat{p}|^2$:
\begin{equation}
T(\hat{p}) = f(|\hat{p}|^2) \label{eq:fhat}
\end{equation}
for some function $f : \mathbb{R}_{\geq 0} \to \mathbb{R}$. \textbf{Anisotropic alternatives dismissed --- Fal-5 dispatched.}

\subsection{C3 forces Taylor expansion with no half-integer powers}

By C3:
\begin{equation}
f(|\hat{p}|^2) = a_0 + a_1 |\hat{p}|^2 + a_2 (|\hat{p}|^2)^2 + a_3 (|\hat{p}|^2)^3 + \ldots
\end{equation}
with $a_0 = 0$ by convention.

A candidate term $a_{1/2} |\hat{p}| = a_{1/2} \sqrt{|\hat{p}|^2}$ is \emph{not analytic} in $|\hat{p}|^2$ at the origin (branch point) and is excluded by C3. The vector linear $\alpha \cdot \hat{p}$ is excluded by C2. \textbf{Falsifier Fal-2 dispatched.}

\subsection{C4 suppresses higher-even-power terms}

In the non-relativistic regime, with the natural dimensional scaling $a_n \sim a_1/(mc)^{2(n-1)}$:
\[
\frac{a_n (|\hat{p}|^2)^n}{a_1 |\hat{p}|^2} \sim \left( \frac{|\hat{p}|^2}{(mc)^2} \right)^{n-1} = \left( \frac{v}{c} \right)^{2(n-1)}.
\]
Higher-order terms are suppressed by powers of $(v/c)^2$ and vanish in the strict $c \to \infty$ limit. Only $a_1 |\hat{p}|^2$ survives. \textbf{Falsifier Fal-3 dispatched.}

\subsection{Relativistic dispersion in non-relativistic regime}

The full relativistic dispersion $T_\mathrm{rel}(\hat{p}) = \sqrt{|\hat{p}|^2 c^2 + m^2 c^4} - m c^2$ Taylor-expands as
\[
T_\mathrm{rel}(\hat{p}) = \frac{|\hat{p}|^2}{2m} - \frac{|\hat{p}|^4}{8 m^3 c^2} + O\!\left( \frac{|\hat{p}|^6}{m^5 c^4} \right).
\]
In the strict non-relativistic limit, only the leading $|\hat{p}|^2/(2m)$ term survives. \textbf{Falsifier Fal-7 dispatched.}

\begin{theorem}[F3]
In the non-relativistic single-particle gauge-coupling-free scope, the kinetic-energy operator $T(\hat{p}) = f(|\hat{p}|^2)$ on the U2 Hilbert space takes the unique form
\[
T(\hat{p}) = a_1 |\hat{p}|^2
\]
with $a_1 > 0$ a real constant determined by F4. All candidate alternatives --- anisotropic, linear-in-$|\hat{p}|$, higher-even-power, full relativistic --- are excluded by C1 + C2 + C3 + C4.
\end{theorem}

\textbf{Falsifiers Fal-2, Fal-3, Fal-5, Fal-7 dispatched.}

\section{Derivation of F4 (Coefficient $1/(2m)$)}\label{sec:F4}

\subsection{The factor-of-2 question}

Dimensional analysis fixes $a_1 = c_1 \hbar^2 / m$ with $c_1$ dimensionless. The standard non-relativistic value is $c_1 = 1/2$, but dimensional analysis alone does not fix it. \textbf{The factor of $1/2$ is the load-bearing question} --- to be derived from primitive-level inputs, not borrowed by convention from classical mechanics.

\subsection{Galilean Lie algebra}

The Galilean Lie algebra is generated by spatial translations $\hat{P}_i$, time translation $\hat{H}$, rotations $\hat{J}_i$, and Galilean boosts $\hat{K}_i$. The structurally substantive commutators are:
\begin{align}
[\hat{K}_i, \hat{P}_j] &= i \hbar m \, \delta_{ij}, \label{eq:KP}\\
[\hat{H}, \hat{K}_i] &= -i \hbar \hat{p}_i. \label{eq:HK}
\end{align}
The mass $m$ enters as the central-extension structure constant. The boost generator takes the form
\begin{equation}
\hat{K}_i = m \hat{x}_i - t \hat{p}_i. \label{eq:K}
\end{equation}

\subsection{Galilean uniqueness within non-relativistic scope}

Within the non-relativistic scope (C4 + absolute time), the Galilean group is the unique consistent boost-translation algebra. Lorentzian boosts mix space and time (relativistic regime); Galilean boosts preserve absolute time (non-relativistic regime). The non-relativistic-limit scope condition selects Galilean over Lorentzian.

\subsection{Galilean transformations act unitarily on the U2 Hilbert space}

\begin{itemize}
  \item Translations $\hat{p}$: unitary by F1 (Stone's theorem).
  \item Time translations $\hat{H}$: assumed self-adjoint under Framing 1.
  \item Rotations $\hat{J}$: unitary by C2.
  \item Boosts $\hat{K} = m \hat{x} - t \hat{p}$: self-adjoint as real-coefficient linear combination of self-adjoint operators.
\end{itemize}
All Galilean transformations are unitary on the U2 Hilbert space.

\subsection{Momentum transformation under boosts}

Under an infinitesimal Galilean boost with velocity $v$:
\[
\hat{p}_j' = \hat{p}_j + \frac{i v_i}{\hbar} [\hat{K}_i, \hat{p}_j] + O(v^2) = \hat{p}_j - m v_j + O(v^2).
\]
For a finite boost:
\begin{equation}
\hat{p} \to \hat{p} - m v. \label{eq:p-boost}
\end{equation}

\subsection{The integration that fixes the factor of 2}

The Hamiltonian $\hat{H} = T(\hat{p})$ must satisfy~\eqref{eq:HK}. Substituting~\eqref{eq:K} and using $[\hat{H}, \hat{p}_i] = 0$:
\[
[\hat{H}, \hat{K}_i] = m [T(\hat{p}), \hat{x}_i] = -i \hbar m \, \frac{\partial T(\hat{p})}{\partial \hat{p}_i}.
\]
Setting equal to $-i\hbar \hat{p}_i$:
\begin{equation}
\frac{\partial T(\hat{p})}{\partial \hat{p}_i} = \frac{\hat{p}_i}{m}. \label{eq:Tprime}
\end{equation}

Using the F3 result $T = f(|\hat{p}|^2)$, so $\partial T / \partial \hat{p}_i = 2 \hat{p}_i f'(|\hat{p}|^2)$:
\[
2 \hat{p}_i f'(|\hat{p}|^2) = \frac{\hat{p}_i}{m} \quad \Longrightarrow \quad f'(|\hat{p}|^2) = \frac{1}{2m}.
\]
Integrating with $f(0) = 0$:
\begin{equation}
f(|\hat{p}|^2) = \frac{|\hat{p}|^2}{2m}. \label{eq:fresult}
\end{equation}

\textbf{The factor of 2 emerges from the integration.} The coefficient $1/(2m)$ in $T = |\hat{p}|^2/(2m)$ comes from integrating $f'(|\hat{p}|^2) = 1/(2m)$ --- the factor of 2 is the chain-rule \emph{Jacobian} relating $\partial T/\partial \hat{p}_i = 2 \hat{p}_i f'(|\hat{p}|^2)$ to the Galilean commutator condition. \textbf{It is not a convention; it is the result of integrating the differential equation that Galilean invariance demands.}

\subsection{Form-FORCED, value-INHERITED}

Combining F3 and F4, restoring $\hbar$:
\begin{equation}
T(\hat{p}) = \frac{\hbar^2 |\hat{p}|^2}{2m} \label{eq:Tfinal}
\end{equation}
with $m$ inherited per Arc~M and $\hbar$ inherited per the dimensional-atlas Madelung anchoring. \textbf{Falsifier Fal-4 dispatched.}

\begin{theorem}[F4]
Under Framing~1 and within the non-relativistic single-particle scope, the kinetic-energy operator $T(\hat{p}) = a_1 |\hat{p}|^2$ established by F3 has
\[
a_1 = \frac{\hbar^2}{2m}
\]
uniquely. The factor of 2 emerges from integrating the Galilean commutator condition $[\hat{H}, \hat{K}_i] = -i\hbar \hat{p}_i$ against $\hat{K} = m\hat{x} - t\hat{p}$, with the factor of 2 the chain-rule Jacobian. The mass $m$ is inherited per Arc~M; $\hbar$ is inherited via the dimensional-atlas Madelung anchoring.
\end{theorem}

\section{Falsifier Resolution Table}\label{sec:falsifier-table}

\begin{center}
\begin{tabular}{p{1.5cm}p{1.0cm}p{8.5cm}p{2.5cm}}
\hline
Falsifier & Target & Dispatch argument & Memo \\
\hline
Fal-1 & F1 & Stone's theorem on U2 Hilbert space's translation group produces unique self-adjoint $\hat{p}$ & Memo 02 \S 1 \\
Fal-2 & F3 & Vector $\alpha \cdot \hat{p}$ excluded by C2; scalar $|\hat{p}| c$ excluded by C3 (branch-point) + C4 & Memo 03 \S 2.4 \\
Fal-3 & F3 & C4 non-rel limit suppresses all $a_n |\hat{p}|^{2n}$ for $n \geq 2$ by $(v/c)^{2(n-1)}$ & Memo 03 \S 2.5 \\
Fal-4 & F4 & Galilean commutator integration: $f'(|\hat{p}|^2) = 1/(2m)$; factor of 2 is integration Jacobian & Memo 03 \S 3.7 \\
Fal-5 & F3 & C2 forces $T(\hat{p}) = f(|\hat{p}|^2)$ & Memo 03 \S 2.3 \\
Fal-6 & F5 & Translation invariance + locality forbid cross terms (gauge-coupling-free scope) & Memo 02 \S 3 \\
Fal-7 & F3 & C4 scope: relativistic dispersion reduces to U4 form in $c \to \infty$ limit & Memo 03 \S 2.6 \\
\hline
\end{tabular}
\end{center}

\textbf{All seven falsifiers dispatched within the stated scope conditions.}

\section{Conditionality Ledger}\label{sec:conditionality}

The U4 verdict is \textsc{forced}, conditional on:

\begin{enumerate}
  \item \textbf{U3 (sibling \textsc{candidate}; Framing 1).} U3 supplies the existence of a Hermitian Hamiltonian generator. U4 specifies the form. Promoting U3 makes U4 fully unconditional.
  \item \textbf{Non-relativistic single-particle scope.} Inherited from Phase-1 Step 2; conditions F3 and F4.
  \item \textbf{Gauge-coupling-free scope.} F5's clean decomposition; gauge field theory is downstream content.
  \item \textbf{Inherited values.} Mass $m$ per Arc~M; $\hbar$ per dimensional-atlas Madelung.
\end{enumerate}

\textbf{U4 introduces no new structural \textsc{candidate}s.} All four conditionalities are inherited from prior structural commitments.

\section{Final U4 Theorem}\label{sec:final-theorem}

\begin{theorem}[U4; Specific Hamiltonian Form]
Let the ED primitive stack supply: bandwidth (Primitive~04), spatial structure (Primitive~06), participation relations (Primitive~03), channel index (Primitive~07), tension polarity (Primitive~09), and relational timing (Primitive~13), together with the now-\textsc{forced} upstream items U1, U2-Discrete, U2-Continuum, U5, and Theorem~10. Within the non-relativistic single-particle scope of Phase-1 Step 2 (gauge-coupling-free), and conditional on the existence of a Hermitian Hamiltonian generator (sibling \textsc{candidate} U3, Framing~1), the Hamiltonian operator on the participation-measure Hilbert space takes the unique form
\[
\hat{H} = \frac{\hbar^2 |\hat{p}|^2}{2m} + V(\hat{x})
\]
with the following structural properties:
\begin{enumerate}
  \item \textbf{Channel--momentum identification.} In the thin/continuum limit, $K \leftrightarrow k$ via Stone's theorem.
  \item \textbf{Kinetic + potential decomposition.} $\hat{H} = T(\hat{p}) + V(\hat{x})$; cross terms forbidden.
  \item \textbf{Quadratic-in-$|\hat{p}|^2$ kinetic structure.} $T(\hat{p}) = a_1 |\hat{p}|^2$ uniquely.
  \item \textbf{Specific coefficient $\hbar^2/(2m)$.} $a_1 = \hbar^2/(2m)$, with the factor of 2 forced by Galilean commutator integration. $m$ inherited per Arc~M; $\hbar$ inherited per dimensional-atlas Madelung.
  \item \textbf{Form $V(\hat{x})$.} Scalar function of position; specific $V$ inherited from external context.
\end{enumerate}
\end{theorem}

\textbf{Status:} \textsc{forced} with form-\textsc{forced}-value-\textsc{inherited} framing.

\section{Downstream Consequences}\label{sec:downstream}

\subsection{Updated structural-foundations ledger}

\begin{center}
\begin{tabular}{cll}
\hline
\# & Theorem & Status \\
\hline
10 & Born rule (Gleason--Busch) & \textsc{forced} \\
11 & U2-Discrete & \textsc{forced} \\
12 & U2-Continuum & \textsc{forced} (with conformal gauge) \\
13 & U5 & \textsc{forced} \\
14 & U1 & \textsc{forced} \\
15 & U4 (this paper) & \textsc{forced} (form; values inherited) \\
\hline
\end{tabular}
\end{center}

\subsection{Active CANDIDATE inventory}

\begin{center}
\begin{tabular}{ll}
\hline
Pre-U4 & Post-U4 \\
\hline
$\{$U3, U4$\}$ + continuum gauge & $\{$U3$\}$ + continuum gauge \\
\hline
\end{tabular}
\end{center}

One upstream \textsc{candidate} remains. Promoting U3 closes the QM-emergence Phase-1 program except for the description-level continuum gauge convention.

\subsection{Schr\"odinger emergence}

The Phase-1 Step 2 chain reads:
\begin{verbatim}
Primitives 03, 04, 06, 07, 09, 13
              |
              v
        U1 (FORCED) -> P_K = sqrt(b) * exp(i pi)
              |
              v
        U2 (FORCED) -> <P|Q> inner product
              |
              v
        U5 (FORCED) -> adjacency partition b = b_x + b_p
              |
              v
        U4 (FORCED, this paper) -> H = h^2 |p|^2/(2m) + V(x)
              |
              v
        U3 (CANDIDATE) -> i h dt P = H P
              |
              v
      Schrodinger emergence (Step 2)
\end{verbatim}

\textbf{Schr\"odinger evolution is \textsc{forced} conditional on U3 alone.}

\subsection{Preparation for the U3 arc}

The U3 arc would address: existence of self-adjoint $\hat{H}$ via Stone's theorem on time-translation symmetry (Primitive~13); linearity; first-order-in-time character. The U3 arc inherits all six \textsc{forced} upstream items. Anticipated arc structure: 4 memos parallel to U4.

\section{Discussion and Outlook}\label{sec:discussion}

\subsection{Position in the Phase-1 program}

U4 is the sixth and most delicate arc of the structural-foundations cycle. The cycle (born\_gleason $\to$ U2-Discrete $\to$ U2-Continuum $\to$ U5 $\to$ U1 $\to$ U4) has demonstrated the structural-derivation methodology against six substantively different structural questions, including the dynamical-content question handled here.

\subsection{Why the factor-of-2 result matters}

The factor of 2 in $1/(2m)$ has appeared in classical mechanics for centuries. In standard quantum mechanics, it is borrowed from the classical formula via $p \to \hat{p}$. The connection is left as a postulate.

The U4 result establishes that this connection is structurally derivable. The factor of 2 is forced by integration of the Galilean Lie algebra commutator condition --- the integration Jacobian of the differential equation $\partial T/\partial \hat{p}_i = \hat{p}_i / m$ produced by the Galilean commutator condition. The factor is not a feature of classical mechanics that quantum mechanics happens to inherit; it is a feature of \emph{Galilean kinematics} that both classical and quantum mechanics inherit.

\subsection{How U4 + U5 + U2 form the dynamical backbone}

The structural-foundations cycle decomposes into three layers:
\begin{itemize}
  \item \textbf{Layer 1 (carrier):} U1.
  \item \textbf{Layer 2 (Hilbert-space arena):} U2-Discrete + U2-Continuum + Theorem~10 + U5.
  \item \textbf{Layer 3 (dynamical):} U3 + U4.
\end{itemize}

U2 + U5 + U4 supply the kinematic + dynamical structural backbone of the Schr\"odinger equation. With U3 closed in the next arc, the full dynamical equation would be structurally complete.

\subsection{Next steps}

\begin{enumerate}
  \item Open the U3 arc.
  \item Revise the QM-emergence Phase-1 synthesis paper.
  \item Draft the U5 publication paper (the missing fourth in the structural-foundations publication series).
\end{enumerate}

\begin{thebibliography}{99}

\bibitem{phase1}
A.~Proxmire, \emph{Phase-1: QM Emergence Structural Completion}, \texttt{papers/QM\_Emergence\_Structural\_Completion/}, 2026.

\bibitem{borngleason}
A.~Proxmire, \emph{The Born Rule as a Forced Theorem of Event Density: A Gleason--Busch Reconstruction from First Principles}, \texttt{papers/Born\_Gleason/}, 2026.

\bibitem{u2}
A.~Proxmire, \emph{The Inner Product as Forced Structure in Event Density: Discrete Derivation, Continuum Lift, and Gauge-Invariant Completion}, \texttt{papers/U2\_Inner\_Product/}, 2026.

\bibitem{u5}
A.~Proxmire, \emph{U5 arc on the adjacency-band Fourier-conjugate partition}, \texttt{arcs/U5/}, 2026.

\bibitem{u1}
A.~Proxmire, \emph{The U1 Theorem: Deriving the Participation Measure from Primitive Structure}, \texttt{papers/U1\_Participation\_Measure/}, 2026.

\bibitem{stone}
M.~H.~Stone, ``On one-parameter unitary groups in Hilbert space,'' \emph{Annals of Mathematics} \textbf{33}, 643--648 (1932).

\bibitem{bargmann}
V.~Bargmann, ``On unitary ray representations of continuous groups,'' \emph{Annals of Mathematics} \textbf{59}, 1--46 (1954).

\bibitem{arcm}
A.~Proxmire, \emph{Arc M: Chain-Mass Structure in Event Density}, \texttt{papers/Arc\_M/}, 2026.

\bibitem{arcn}
A.~Proxmire, \emph{Arc N: Non-Markovian Structure as a Forced Memory-Kernel Layer of Event-Density Theory}, \texttt{papers/Arc\_N/}, 2026.

\bibitem{arcr}
A.~Proxmire, \emph{Arc R: Relativistic Kinematics and Dynamics in Event Density}, \texttt{papers/Arc\_R/}, 2026.

\end{thebibliography}

\vspace{1em}
\noindent\rule{\linewidth}{0.4pt}
\vspace{0.2em}

\noindent\emph{Manuscript prepared April 2026. Source materials: \texttt{arcs/U4/} four-memo arc and closure memo, integrated into a single publishable document. The U4 arc is the sixth and most delicate of the structural-foundations cycle.}

\noindent\emph{Draft complete. Pausing for revision.}

\end{document}
